\section{Source objects and the finite calculation}
This calculation uses the preserved 165-page supplied manuscript
\emph{Finite Time Blowup for Navier--Stokes}, source equations
(7.23)--(7.42), (9.12)--(9.14). Its PDF SHA-256 is
\begin{center}\small
\texttt{8c8a94ad9ac824c8b605b9827cadf7beaca48bd10b380de3cfc872a2c37afa81}.
\end{center}
The source pulse fields, auxiliary rectangles, phases, and leading
covariance columns are the given objects; their global construction is
source input. Every finite map asserted below is proved here. The existing
coupled-stage proofs remain in the complete 102-page reader. The companion
pulse note proves the exact frame residual and a constructed finite-interval
intertwiner along its evolving coefficients. The source uses positive
$q(z,t)$, band scales $Q_\beta=2^{-\ell_\beta}$,
$\varepsilon_\beta=Q_\beta^h$, $A=1/2+h$; none is set to one here or
identified with the original coupled $A_0,\lambda_0,\mu,\nu=\sqrt{A_0}$
or its physical diffusivities.

All operations in this note are on a compact set inside $q>0$ and the
source active annulus. There are finitely many relevant bands and
locally finite source labels. Supports of newly prescribed inputs are
strictly inside the annulus, so no lower bound on a leading amplitude
at a vanishing shell edge is assumed. Endpoint and shell-edge estimates
are not inferred from compact-interior estimates.

\section{Physical curl of the actual retained amplitude}
Let $y=(\Theta,\Omega)^T$ be the actual smooth retained coupled solution,
$y'=q_v(C_cy+c_c)$, with the coefficients of the companion pulse note.
Use the smooth actual coefficients on the chosen slow/auxiliary patch;
all harmonic coefficients and cutoffs are independent of $\theta$ in
cylindrical components, and each carrier has nonzero integer angular
frequency. For its source tangent frame $B$, smooth coefficient map $L$,
and a fixed smooth real compact cutoff $\chi$ of this type, put
\[
 t_m=\chi BLy,\qquad
 \mathcal E_\chi=
 \chi\{[L'+q_vLC_c-HL+m^2d_NL]y+q_vLc_c\}+\chi'Ly,
 \quad H=B_\ell(\mathsf A_\Phi B-B').
\]
Here prime differentiates along the pulse coordinate; the other slow
and auxiliary derivatives remain in the spatial formulas below.
Smooth parameter dependence of the companion's constructed $L$ follows
directly from its integral equations on a fixed interval: for a nonzero
parameter multi-index $I$,
\[
 \partial^I X(v)=
 \int_{v_0}^v\left[D\,\partial^I X+
 \sum_{0<J\le I}{I\choose J}(\partial^JD)
                                 \partial^{I-J}X\right]\,\dd w .
\]
The initial parameter derivatives vanish for identity initial data.
Inductively the known forcing is smooth and bounded on the compact patch;
the same convergent ordered-integral iteration solves this equation.
This proves every finite mixed derivative of $X,X^{-1},L$.
Since $n^TB=0$, $n^Tt_m=0$. Differentiation proves
$t_m'-\mathsf A_\Phi t_m+m^2d_Nt_m=B\mathcal E_\chi$.
Thus even the pulse-cutoff term is explicit.

A physical version of the exact curl construction makes the intermediate
map to momentum residual unambiguous. For each chosen harmonic let
$\phi(x,t)=k_\beta m\Phi_\beta(x,t,Y(x,t))$ be its real evaluated phase,
let $\xi=\nabla_x\phi\ne0$, and let $a(x,t)=Q_\beta^{-A}t_m$ be the
physical transverse amplitude. The source coordinate change is
$r=\sqrt{Q_\beta}R$, $z=Q_\beta^D Z$, $1-t=Q_\beta T$,
with $D=1/2-h$ and normalized axial derivative
$D_z=\varepsilon_\beta\partial_Z$. It gives
$\xi=(k_\beta m/\sqrt{Q_\beta})n_\Phi$; all auxiliary chain derivatives
are included in $n_\Phi$. Hence $\xi\cdot a=0$. Define in Cartesian
components
\begin{equation}
 C=\frac{i\,\xi\times a}{|\xi|^2},\quad
 \mathcal A=C e^{i\phi},\quad
 \nabla\times\mathcal A=(a+r_a)e^{i\phi},\qquad
 r_a=\nabla\times C. \label{cb:curl}
\end{equation}
The vector triple-product identity gives
$i\xi\times C=a$, proving the last equality. It also shows that the
physical potential has exactly the source factor $Q_\beta^{1/2-A}$
relative to its chart coefficient
$i\,n_\Phi\times t_m/(k_\beta m|n_\Phi|^2)$.
Cartesian differentiation of $C$ includes derivatives of the cylindrical
basis; the same terms are the $R^{-1}$ terms of source (7.37)--(7.38).
Take real parts and sum the finite retained harmonics. This constructs
a smooth compact velocity $w_y$ with $\nabla\cdot w_y=0$.

Retain a finite smooth source state $(u_*,p_*)$, including its background,
actual leading waves and their curl corrections, and the real
physical pressure increments obtained from the companion pulse pressure
formula with $f_m=0$:
\[
 \pi_{\chi,m}=\frac{i}{k_\beta m|n|^2}
                  \bigl(n^T\mathsf Kt_m-(n')^Tt_m\bigr),\qquad
 \pi_y=\sum_{\beta,m}\operatorname{Re}
             (Q_\beta^{-2A}\pi_{\chi,m}e^{i\phi_{\beta,m}}).
\]
These are smooth and compact wherever $t_m$ is. This definition fixes the
pressure map without assuming the pulse residual vanishes.
Define the actual candidate residual
\begin{align}
 R_y={}&R_{\nu_{\rm NS}}(u_*+w_y,p_*+\pi_y), \nonumber\\
 R_{\nu_{\rm NS}}(u,p)
   :={}&\partial_tu+(u\cdot\nabla)u-\nu_{\rm NS}\Delta u+\nabla p .
 \label{cb:Ry}
\end{align}
This is a definition by explicit differentiated fields, not an assumed
source term. Its exact expansion is
\[
 R_y=R_{\nu_{\rm NS}}(u_*,p_*)
 +\partial_tw_y+(u_*\cdot\nabla)w_y+(w_y\cdot\nabla)u_*
 -\nu_{\rm NS}\Delta w_y+\nabla\pi_y
 +(w_y\cdot\nabla)w_y .
\]
Every cutoff, parent, phase, and curl derivative occurs in this formula.
Physical $\nu_{\rm NS}$ is retained and is distinct from every source
time-scaling convention. Define
$\Delta R_y=R_y-R_{\nu_{\rm NS}}(u_*,p_*)$.
We correct the new compact increment; this difference has compact support
since every added term does. The baseline residual remains explicitly.

The averaging needed next is the source extended-field operation, not a
spatial average of the evaluated field. Evaluate derivatives on the
extended coefficients first by the complete chain rule. In a fixed
chart set $\mathcal D_i=\partial_{x_i}
+\sum_\alpha(\partial_{x_i}Y_\alpha)\partial_{Y_\alpha}$ and similarly
$\mathcal D_t$. These are the pullbacks of commuting physical
derivatives, so $[\mathcal D_i,\mathcal D_j]=0$ on their lifted domain.
For smooth periodic coefficients their auxiliary derivative terms have
zero Haar average. In cylindrical components angular differentiation
also differentiates the basis. This defines an extended residual
$\widetilde R_y$ whose evaluation is exactly \eqref{cb:Ry}.
Average the tangential components of the increment at fixed slow variables:
\begin{equation}
 F_2=\av{\Delta\widetilde R_{y,\theta}},\qquad
 F_1=\av{\Delta\widetilde R_{y,z}},\qquad
 \av{g}=\int_{\mathbb T^2}\frac1{2\pi}
                    \int_0^{2\pi}g\,\dd\theta\,\dd Y . \label{cb:F}
\end{equation}
These $F_e$ are smooth and compact
in the chosen physical annulus. Equations \eqref{cb:curl}--\eqref{cb:F}
are the explicit intermediate map from the retained coupled pair to a
signed momentum residual. They do not identify it with the original
two-dimensional Boussinesq force.

\section{Radial moments, support, and the original chart scales}
For each fixed $(z,t)$, choose $0<r_-<r_+$ containing the compact
supports of both $F_e$ and the chosen bumps in the active annulus.
The source bump is
\[
 b_e(r,z,t)=q^{-(e+1)/2}\widehat b_e(r/\sqrt q),\qquad
 \int_0^\infty x^e\widehat b_e(x)\,\dd x=1,\qquad e=1,2,
\]
with interior support. The change $r=\sqrt q\,x$ proves
$\int r^eb_e\,\dd r=1$ exactly. Define
\begin{equation}
 M_e=\int_0^\infty r^eF_e(r)\,\dd r,\qquad
 \sigma_e(r)=-r^{-e}\int_0^r s^e(F_e(s)-b_e(s)M_e)\,\dd s .
 \label{cb:primitive}
\end{equation}
Then
\begin{equation}
 (\partial_r+e/r)\sigma_e=-F_e+b_eM_e,\qquad
 \int_0^\infty r^e(F_e-b_eM_e)\,\dd r=0. \label{cb:radial}
\end{equation}
Differentiate $r^e\sigma_e$ to prove the first identity; integrate and
use the bump normalization for the second. The integral is zero below
$r_-$ and above $r_+$, so $\sigma_e$ has compact smooth support in the
same enclosing interval. A homogeneous solution is $cr^{-e}$, and a
compact homogeneous solution has $c=0$; thus the primitive is unique
among compact solutions of its displayed equation. Moreover a compact
solution of $(\partial_r+e/r)\sigma=-F_e$ exists if and only if $M_e=0$:
necessity follows by integrating the weighted derivative, and
sufficiency follows from \eqref{cb:primitive}. The two moments are
therefore the exact surviving finite-dimensional terms.

For $C_e=\int_{r_-}^{r_+}s^e\,\dd s$,
$|M_e|\le C_e\|F_e\|_\infty$ and
\[
 \|\sigma_e\|_\infty\le r_-^{-e}C_e
        (\|F_e\|_\infty+\|b_e\|_\infty|M_e|).
\]
If $G_e=F_e-b_eM_e$, successive radial derivatives obey the exact recurrence
\[
 \partial_r^{j+1}\sigma_e=-\partial_r^jG_e
 -e\sum_{i=0}^j {j\choose i}(-1)^ii!r^{-i-1}
                              \partial_r^{j-i}\sigma_e .
\]
This proves every finite radial derivative bound inductively. Parameter
derivatives follow by differentiating the compact integral; all derivatives
of $q,b_e$ and the moving physical data are retained. On the present compact
set $q$ is bounded away from zero. No terminal uniform estimate is implied.

The source (9.12)--(9.14) uses physical
$F_e=Q^{-2A-1/2}E_e^*$, where
$E_2^*=\av{E_\theta}_Y$ and $E_1^*=\av{E_z}_Y$.
For our fields define $E_e^*=Q^{2A+1/2}F_e$ in that same convention.
With $r=\sqrt Q R$ and $Q$ fixed in a chart, the exact conversions are
\begin{gather}
 M_e=Q^{e/2-2A}M_e^*,\quad M_e^*=\int R^eE_e^*\,\dd R,\quad
 b_e^*(R)=Q^{(e+1)/2}b_e(\sqrt Q R),\nonumber\\
 \Sigma_e(R)=Q^{2A}\sigma_e(\sqrt Q R),\qquad
 (\partial_R+e/R)\Sigma_e=-E_e^*+b_e^*M_e^*. \label{cb:chart}
\end{gather}
Substitute $r=\sqrt Q R$ including its measure in the moment, and
differentiate $\Sigma_e$ to prove each factor. The input $F_e$ is already
auxiliary independent. Taking only an averaged moment of an unaveraged
$F(R,Y)$ would not suffice: $F=b_e(R)\cos(2\pi Y_1)$ has zero averaged
moment but a nonzero exterior primitive at each generic $Y$.

\section{The source covariance derivative and its signed right inverse}
At each slow point let $w_{\rm tan}=(w_\theta,w_z)$ and define
\[
 \Cov(w)=\av{w_rw_{\rm tan}},\qquad
 \Cross(w,v)=\av{w_rv_{\rm tan}+v_rw_{\rm tan}}.
\]
Direct expansion proves
$\Cov(w+v)=\Cov(w)+\Cross(w,v)+\Cov(v)$ and
$D\Cov(w)[v]=\Cross(w,v)$.

Use the exact source real pulses
$b_\sigma=\chi_g(\xi_g)\psi(v)t_\sigma^h\cos(k\Phi_\sigma)$,
$\sigma\in\{+,-\}$. They exclude the slow cutoff $\eta_\beta$ and
scalar amplitudes; $\chi_g$ is an auxiliary transverse cutoff.
The two signs have disjoint auxiliary supports. Put
$H_{\rm cov}=[\Cov(b_+)\mid\Cov(b_-)]$.
For slow scalar coefficients independent of both averages,
\[
 \Cov(a_+b_++a_-b_-)=H_{\rm cov}(a_+^2,a_-^2)^T.
\]
All cross products vanish by disjointness. The factor $1/2$ from
$\av{\cos^2(k\Phi_\sigma)}_\theta$ is already inside the columns.

On the source interior cone the fixed primary amplitudes satisfy
$a_\sigma=\sqrt{(H_{\rm cov}^{-1}T_{0,*})_\sigma}>0$,
$T_{0,*}=(Q/q)^{A+1/2}T_0$, and
$W_0=\sqrt\varepsilon\sum_\sigma a_\sigma b_\sigma$.
For any signed slow real vector $\Sigma$, define
\begin{equation}
 d_\Sigma=H_{\rm cov}^{-1}\Sigma/\varepsilon,\qquad
 \delta a_\sigma=(d_\Sigma)_\sigma/(2a_\sigma),\qquad
 L_\Sigma=\sqrt\varepsilon\sum_\sigma\delta a_\sigma b_\sigma .
 \label{cb:L}
\end{equation}
It follows exactly that
\begin{gather}
 \Cross(W_0,L_\Sigma)=\varepsilon H_{\rm cov}(2a\odot\delta a)=\Sigma,
 \label{cb:inverse}\\
 \Cov(W_0+L_\Sigma)
   =\varepsilon T_{0,*}+\Sigma+
       \varepsilon H_{\rm cov}(\delta a_+^2,\delta a_-^2)^T.
 \label{cb:quad}
\end{gather}
This proves a linear right inverse of the derivative. It is not an
exact inverse of the quadratic covariance map. There is no sign or
size restriction on $\Sigma$ in these identities; no square root of
$H_{\rm cov}^{-1}(T_{0,*}+\Sigma/\varepsilon)$ is taken.

With Euclidean and induced operator norms and
$a_{\min}=\min_\sigma|a_\sigma|>0$, they give the explicit bound
\begin{equation}
 \|\Cov(L_\Sigma)\|\le
 \frac{\|H_{\rm cov}\|\|H_{\rm cov}^{-1}\|^2}
              {4\varepsilon a_{\min}^2}\,\|\Sigma\|^2. \label{cb:bound}
\end{equation}
Indeed $\|\delta a\|\le\|H_{\rm cov}^{-1}\|\|\Sigma\|/
(2\varepsilon a_{\min})$ and
$\|(\delta a_\sigma^2)_\sigma\|_2\le\|\delta a\|_2^2$.
All higher derivative costs are explicit from Leibniz and
\[
 \partial^I H_{\rm cov}^{-1}
 =-H_{\rm cov}^{-1}
 \sum_{0<J\le I}{I\choose J}(\partial^JH_{\rm cov})
                           \partial^{I-J}H_{\rm cov}^{-1},
\]
and, for any nonzero scalar $a$,
$\partial^I(a^{-1})=-a^{-1}\sum_{0<J\le I}{I\choose J}
(\partial^Ja)\partial^{I-J}(a^{-1})$.
Together with \eqref{cb:L} these recurrences give finite bounds for
each order without assuming constant amplitudes or constant columns.

\section{Assembly, the actual covariance change, and residual cancellation}
For the physical target $\sigma=(\sigma_2,\sigma_1)$ from
\eqref{cb:primitive}, use the source squared slow partition
$\sum_\beta\eta_\beta^2=1$ on its support. Set
\begin{align}
 W_{0,\rm phys}&=\sum_\beta Q_\beta^{-A}\eta_\beta W_{0,\beta},\nonumber\\
 v_{\rm pr}&=\sum_\beta Q_\beta^{-A}\eta_\beta
                        L_\beta(Q_\beta^{2A}\sigma).
 \label{cb:assembly}
\end{align}
Overlapping slow labels have disjoint auxiliary supports. Taking cross
covariance and using \eqref{cb:inverse} gives
\[
 \Cross(W_{0,\rm phys},v_{\rm pr})
 =\sum_\beta Q_\beta^{-2A}\eta_\beta^2Q_\beta^{2A}\sigma=\sigma.
\]
Every band scale remains distinct. Likewise the leading physical factor is
$Q^{-2A}\varepsilon(Q/q)^{A+1/2}
=Q^{-A+h+1/2}q^{-A-1/2}=q^{-A-1/2}$ because $A=1/2+h$.

Apply the exact curl \eqref{cb:curl} to each coefficient of
$v_{\rm pr}$, including the slow cutoffs before curling. Write the
resulting divergence-free increment as $v_s=v_{\rm pr}+r_s$.
Let $t_{s,m}$ denote the chart transverse harmonic coefficients of
$\eta_\beta L_\beta(Q_\beta^{2A}\sigma)$, before multiplication by
$Q_\beta^{-A}$. For each such coefficient choose
$\pi_{s,m}=i(n^T\mathsf Kt_{s,m}-(n')^Tt_{s,m})/
(k_\beta m|n|^2)$ and sum
$\operatorname{Re}(Q_\beta^{-2A}\pi_{s,m}e^{i\phi_{\beta,m}})$
to define the smooth compact $\pi_s$. Its full gradients remain below.
Let the actual old wave $w$ consist of the wave part of $u_*$
plus the candidate $w_y$, including all existing curl corrections, and write
$w=W_{0,\rm phys}+e_w$. Define tensor covariance
$\mathcal T(w)=\av{w\otimes w}$ and
$\mathcal B_T(u,v)=\av{u\otimes v+v\otimes u}$.
Expansion, with $e_w=w-W_{0,\rm phys}$, proves
\begin{equation}
 \Delta\mathcal T=
 \mathcal B_T(W_{0,\rm phys},v_{\rm pr})
 +\mathcal B_T(e_w,v_{\rm pr})
 +\mathcal B_T(w,r_s)+\mathcal T(v_s).
 \label{cb:fullcov}
\end{equation}
Its first radial--tangential pair is $\sigma$.
All other entries and the remaining three tensors are retained.
For example their radial--tangential pair has bound
\[
 2\|e_w\|_{L^2_{\theta,Y}}\|v_{\rm pr}\|_{L^2_{\theta,Y}}
 +2\|w\|_{L^2_{\theta,Y}}\|r_s\|_{L^2_{\theta,Y}}
 +\|v_s\|_{L^2_{\theta,Y}}^2 ,
\]
by Cauchy--Schwarz. The compact-set estimates above and the explicit
curl differentiate to bound all these terms at each finite order.

For clarity the complete residual after adding $v_s$ and its selected
pressure $\pi_s$ is, without any averaging,
\begin{align}
 R_{\nu_{\rm NS}}(u+v_s,p+\pi_s)
 =R_{\nu_{\rm NS}}(u,p)
 &+\partial_tv_s+(u\cdot\nabla)v_s+(v_s\cdot\nabla)u \nonumber\\
 &-\nu_{\rm NS}\Delta v_s+\nabla\pi_s
 +(v_s\cdot\nabla)v_s . \label{cb:residual}
\end{align}
This follows by expanding the product and retains viscosity and every
pressure term. For a symmetric tensor $S$ the averaged tangential
divergence in physical cylindrical coordinates is
\[
 (\av{\nabla\cdot S})_\theta
   =(\partial_r+2/r)\av{S_{r\theta}}+\partial_z\av{S_{z\theta}},
 \qquad
 (\av{\nabla\cdot S})_z
   =(\partial_r+1/r)\av{S_{rz}}+\partial_z\av{S_{zz}} .
\]
To verify, write the Cartesian divergence in the orthonormal cylindrical
basis, using $\partial_\theta e_r=e_\theta$,
$\partial_\theta e_\theta=-e_r$; periodic angular derivatives average
to zero and the two basis terms give $2/r$ and $1/r$ respectively.
The auxiliary chain terms average to zero by periodicity, but the
physical slow derivatives remain.

Let $U=u_*+w_y$, $P_y=p_*+\pi_y$, and set
\[
 \mathcal L=
 \partial_tv_s+(U\cdot\nabla)v_s+(v_s\cdot\nabla)U
 -\nu_{\rm NS}\Delta v_s+\nabla\pi_s .
\]
For a fully explicit averaged remainder, write $U=u_{\rm mean}+w$,
where $u_{\rm mean}$ is its angular mean, and put
$\mathcal V=\mathcal L-(w\cdot\nabla)v_s-(v_s\cdot\nabla)w$.
Then the average of \eqref{cb:residual} is
$\av{R_y}+\av{\mathcal V}+\nabla\cdot\Delta\mathcal T$.
Let $S_0=\mathcal B_T(W_{0,\rm phys},v_{\rm pr})$ and
$S_{\rm rem}=\mathcal B_T(e_w,v_{\rm pr})
 +\mathcal B_T(w,r_s)+\mathcal T(v_s)$.
Consequently the radial divergence of the first pair in
\eqref{cb:fullcov} changes the prescribed averaged residual $F_e$ into
$b_eM_e$ by \eqref{cb:radial}. The axial tensor divergences, the other
three tensors, the exact linear terms and the pressure terms of
\eqref{cb:residual} remain in the new residual. Equivalently the new
averaged tangential residual is exactly
\begin{align*}
 \av{R_{\rm new}}_{\rm tan}
 ={}&\av{R_{\nu_{\rm NS}}(u_*,p_*)}_{\rm tan}
 +(b_2M_2,b_1M_1)+\av{\mathcal V}_{\rm tan}\\
 &+\partial_z((S_0)_{z\theta},(S_0)_{zz})
   +(\nabla\cdot S_{\rm rem})_{\rm tan}.
\end{align*}
Every term here is defined by the displayed differentiated fields and
tensors. None is assumed to vanish.
The map from $y$ through \eqref{cb:F}, \eqref{cb:primitive},
\eqref{cb:L}, and \eqref{cb:assembly} is now an actual composite
construction. It retains the signed residual produced by the coupled
candidate and the source's original physical and chart units.

\section{Signs and the exact full-flow reflection}
The maps $F_e\mapsto M_e\mapsto\sigma_e\mapsto L_\Sigma$ are linear
and reverse sign with their inputs; $\Cov(L_\Sigma)$ is even.
The two primary amplitudes may also be replaced independently by
$s_\sigma a_\sigma$, $s_\sigma\in\{-1,1\}$, provided the same fixed
choice is used in each denominator: the cross identity is unchanged.
This does not negate the nonlinear Navier--Stokes equation.

Here is the exact full-flow map. In Cartesian coordinates let
$O=\operatorname{diag}(1,-1,1)$ and define
\[
 u^\sharp(x,t)=O u(Ox,t),\quad
 p^\sharp(x,t)=p(Ox,t),\quad f^\sharp(x,t)=O f(Ox,t).
\]
The chain rule gives
$\nabla u^\sharp=O(\nabla u)(Ox,t)O$,
$\Delta u^\sharp=O(\Delta u)(Ox,t)$, and
$(u^\sharp\cdot\nabla)u^\sharp
 =O((u\cdot\nabla)u)(Ox,t)$.
Thus $\nabla\cdot u^\sharp=0$ and
$R_{\nu_{\rm NS}}(u^\sharp,p^\sharp)=O R_{\nu_{\rm NS}}(u,p)(Ox,t)$
with the same positive viscosity.
In cylindrical components this is
\[
 (u_r^\sharp,u_\theta^\sharp,u_z^\sharp)(r,\theta,z,t)
 =(u_r,-u_\theta,u_z)(r,-\theta,z,t),
\]
and the force transforms identically. Pressure reflects in its argument.
For an axisymmetric field this reverses exactly its swirl; meridional
components remain fixed. The entire momentum tensor is mapped to
$O\mathcal T O^T$, so its radial--tangential pair transforms by
$(T_{r\theta},T_{rz})\mapsto(-T_{r\theta},T_{rz})$ after angular
averaging. Vorticity transforms as
$\omega^\sharp(x,t)=\det(O)\,O\omega(Ox,t)$, obtained by applying
the same derivative identities to the Levi-Civita formula for curl.
The map is an involution, preserves the kinetic-energy integral
(by $|\det O|=1$), compact support, smoothness, and blowup magnitude
along the reflected path. It preserves zero initial data when present.
It proves transport of any established axisymmetric signed swirl
divergence to the opposite sign; it does not itself prove an endpoint.

The source five-moment correction, auxiliary mean inverse, and infinite
summation are subsequent source operations. Their full estimates are
not re-proved here. The earlier draft formulas for them are withdrawn:
in particular the fifth moment row is
$\int(2RG\gamma_d-RV\Delta v)\,\dd R=-J_z$, without an extra outer
$R$, and its reserved profile in a $Q$ chart includes $(Q/q)^A$.
The finite map proved above retains the two moments and all other
residual terms, so it establishes neither a third shooting return nor
an infinite modified viscous sequence.
